%==============================================================
% Stikkord til kavitasjons-diskusjon (Ch5 Discussion)
%==============================================================

\subsection{Cavitation margin assessment}
\label{sec:cav_discussion}

%-- Setup-feil i simuleringen ---------------------------------
\paragraph{Boundary condition offset (Methods context).}
\begin{itemize}
    \item Reference Pressure $P_\mathrm{ref} = \SI{1}{atm}$ combined with Inlet Static Pressure of \SI{1}{atm} (relative) gave an effective absolute inlet of $\sim\!\SI{2}{atm}$.
    \item This corresponds to $\sim\!\SI{10.35}{m}$ of artificial suction head, while the intended physical setup has $\sim\!\SI{0.95}{m}$ tank height or, conservatively, \SI{0}{m} suction head.
    \item Consequence for the main GV screening (head loss $\Delta H$, swirl, ranking across the 18 cases): \textbf{none} -- only absolute pressure values are shifted by a constant offset.
    \item Consequence for cavitation: \textbf{must be re-evaluated} using absolute pressures or an equivalent zero-suction-head interpretation.
\end{itemize}

%-- Korrekt formel --------------------------------------------
\paragraph{Cavitation margin formulation.}
The physically correct margin is
\begin{equation}
    \mathrm{cavMargin}_H \;=\; \frac{p_\mathrm{min,abs} - p_\mathrm{vap}}{\rho g},
    \label{eq:cavmargin}
\end{equation}
where $p_\mathrm{min,abs}$ is the minimum \emph{absolute} static pressure on the wetted surfaces and $p_\mathrm{vap} \approx \SI{3170}{Pa}$ at laboratory temperature.
\begin{itemize}
    \item The CFX CEL expression \texttt{(pMin - pVap)/($\rho$g)} uses pressure relative to $P_\mathrm{ref}$.
    \item For $P_\mathrm{ref} = \SI{1}{atm}$, this CEL output corresponds directly to a \textbf{zero-suction-head} interpretation -- the worst-case installation assumption.
    \item This is the assumption adopted for the cavitation analysis in this thesis.
\end{itemize}

%-- Resultater ------------------------------------------------
\paragraph{Cavitation margins for case CC\_250\_9.}
\begin{table}[h]
\centering
\begin{tabular}{lrr}
\hline
Location & $p_\mathrm{min}$ (relative) & $\mathrm{cavMargin}_H$ \\
\hline
GV passage & $+\SI{75.5}{kPa}$ & $+\SI{7.40}{m}$ (safe) \\
RDT blades & $-\SI{19.0}{kPa}$ & $\boldsymbol{-\SI{2.27}{m}}$ \textbf{(cavitating)} \\
\hline
\end{tabular}
\end{table}

%-- Tolkning --------------------------------------------------
\paragraph{Interpretation.}
\begin{itemize}
    \item A negative margin indicates that the local pressure on the RDT blade suction side drops below the vapour pressure -- vapour cavities are predicted to form.
    \item This is consistent with the hubless rim-driven configuration, where high tip-region velocities produce local low-pressure zones not present in conventional hubbed designs.
    \item The GV section retains a comfortable margin ($+\SI{7.4}{m}$); it is \textbf{not} the cavitation-limiting component.
    \item The RDT is the cavitation-critical component in the system.
\end{itemize}

%-- Installasjons-dykk ----------------------------------------
\paragraph{Required submergence for cavitation-neutral operation.}
\begin{itemize}
    \item Closing a margin gap of \SI{2.27}{m} requires \SI{2.27}{m} of additional suction head -- equivalent to submerging the RDT centerline \SI{2.27}{m} below the tank free surface (neglecting suction-line losses).
    \item The general NPSH balance:
    \begin{equation}
        \mathrm{NPSH}_a \;=\; h_\mathrm{atm} + h_\mathrm{suction} - h_\mathrm{vap} - h_\mathrm{loss,suction}
        \label{eq:npsh}
    \end{equation}
    with $h_\mathrm{atm} \approx \SI{10.35}{m}$, $h_\mathrm{vap} \approx \SI{0.32}{m}$.
    \item Minimum submergence in this analysis is set by the most negative blade-surface pressure, not by the bulk inlet condition.
\end{itemize}

%-- Sikkerhetsmargin ------------------------------------------
\paragraph{Safety margins.}
\begin{itemize}
    \item Recommended $25\text{--}50\%$ margin above the minimum to accommodate:
    \begin{itemize}
        \item Steady RANS underestimates local low-pressure peaks (transient, tip-vortex, and blade-passing fluctuations not captured).
        \item Off-design operation: higher flow rates raise local velocities and lower $p_\mathrm{min}$.
        \item Variations in $p_\mathrm{vap}$ with operating temperature.
    \end{itemize}
    \item $25\%$ margin: submergence $\approx \SI{2.84}{m}$.
    \item $50\%$ margin: submergence $\approx \SI{3.4}{m}$.
    \item Conservative recommendation for a prototype: $\gtrsim \SI{3}{m}$.
\end{itemize}

%-- Begrensninger ---------------------------------------------
\paragraph{Limitations and caveats.}
\begin{itemize}
    \item The simulation does \textbf{not} model phase change -- predicted minima are indicative only; actual vapour formation would alter the local flow field.
    \item Steady-state RANS typically underpredicts transient low-pressure events; the true cavitation onset may be more severe.
    \item Hubless rim-driven configurations are particularly sensitive to tip-region cavitation.
    \item The reported margin applies at a single operating point; off-design conditions are not assessed.
\end{itemize}

%-- Kobling til GV-screening ----------------------------------
\paragraph{Coupling to GV design (link to main screening).}
\begin{itemize}
    \item RDT cavitation is governed primarily by RDT blade geometry, not by GV design.
    \item However, GV designs with high passage loss reduce absolute pressure throughout the downstream section, which can influence the coupled RDT--GV--RPT NPSH balance.
    \item Future ranking of the 18 GV cases may include $\mathrm{cavMargin}_\mathrm{RDT}$ as a fourth screening criterion alongside $\Delta H$, $S$, and outlet flow uniformity.
\end{itemize}

%-- Anbefalt formulering --------------------------------------
\paragraph{Summary statement for the thesis.}
The simulation boundary conditions correspond to $\sim\!\SI{2}{atm}$ absolute inlet ($\sim\!\SI{10.35}{m}$ suction head); for cavitation assessment, results are reinterpreted under a zero-suction-head assumption representing a worst-case installation. Under this assumption, the cavitation head margin at the RDT minimum-pressure location is $-\SI{2.27}{m}$, indicating onset of cavitation in the hubless tip region under atmospheric inlet conditions. The GV sections retain a margin of $+\SI{7.4}{m}$ and are not limiting. To achieve cavitation-neutral operation, the RDT centerline must be submerged at least \SI{2.27}{m} below the tank free surface; a conservative installation with a $25\text{--}50\%$ safety margin gives \SIrange{2.8}{3.4}{m}. The RDT is therefore identified as the cavitation-critical component, and cavitation-aware redesign or installation adjustment is recommended for prototype implementation.

%-- Future Work -----------------------------------------------
\paragraph{Future Work links.}
\begin{itemize}
    \item Re-simulate with correct boundary conditions: Inlet relative pressure set to $\rho g h_\mathrm{tank}$ for the chosen prototype submergence.
    \item Transient / URANS to capture blade-passing pressure peaks.
    \item Explicit cavitation modelling (Rayleigh--Plesset or Zwart--Gerber--Belamri) to resolve phase change and its effect on performance.
    \item Experimental verification of cavitation inception across a range of suction heads.
\end{itemize}